\documentclass[a4paper, serif, 10pt]{moderncv}

%% moderncv
% color and style
\moderncvcolor{black}
\moderncvstyle{banking}

%% page layout/margins
\usepackage[top=2.5cm, bottom=2.5cm, left=1.5cm, right=1.5cm]{geometry}

\nopagenumbers{}

%% Babel config for Italian language
% ref:
% - https://latex3.github.io/babel/guides/locale-italian.html
\usepackage[italian]{babel}

%% fontspec
\usepackage{fontspec}

\IfFontExistsTF{Linux Libertine}{
  \setmainfont[Ligatures={TeX, Common}, Numbers=OldStyle]{Linux Libertine}
}{}
\IfFontExistsTF{Linux Libertine O}{
  \setmainfont[Ligatures={TeX, Common}, Numbers=OldStyle]{Linux Libertine O}
}{}

%% CTeX
% CJK support, used to show CJK characters in the resume
%
% - heading=false: disable ctex heading
% - punct=kaiming: use kaiming punctuations styles for CJK
% - scheme=plain: use plain scheme, do not override \normalsize font size
% - space=auto: space settings for CJK characters
%
% CJK fonts are resolved per locale language, e.g. Noto Serif/Sans CJK JP
% for Japanese resumes, with fallback to the next region if not installed
%
% ref:
% - http://ctan.mirrorcatalogs.com/language/chinese/ctex/ctex.pdf
\usepackage[UTF8, heading=false, punct=kaiming, scheme=plain, space=auto]{ctex}

\IfFontExistsTF{Noto Serif CJK SC}{
  \setCJKmainfont{Noto Serif CJK SC}
}{}
\IfFontExistsTF{Noto Sans CJK SC}{
  \setCJKsansfont{Noto Sans CJK SC}
}{}

%% line spacing
\usepackage{setspace}
\setstretch{1.125}

%% URL styling - use normal text instead of monospace
\urlstyle{same}

%% Patch \httplink and \httpslink to handle full URLs
% The original moderncv macros always prepend http:// or https:// to URLs.
% This causes issues when the URL already has a protocol (e.g., https://example.com)
% resulting in malformed URLs like https://https://example.com.
% This patch checks if the URL already contains "://" and uses it directly if so.
% Note: We use \str_set:Nx (x-expansion) to expand macros like \@homepage first.
\ExplSyntaxOn
\str_new:N \g_yamlresume_url_str

\RenewDocumentCommand{\httpslink}{O{}m}{
  \str_set:Nx \g_yamlresume_url_str {#2}
  \str_if_in:NnTF \g_yamlresume_url_str {://}
    { \url{#2} }
    { \url{https://#2} }
}

\RenewDocumentCommand{\httplink}{O{}m}{
  \str_set:Nx \g_yamlresume_url_str {#2}
  \str_if_in:NnTF \g_yamlresume_url_str {://}
    { \url{#2} }
    { \url{http://#2} }
}
\ExplSyntaxOff

\name{Luca Moretti}{}
\title{Ingegnere AI senior specializzato in apprendimento automatico e MLOps}
\phone[mobile]{+39 347 892 1456}
\email{luca.moretti.ai@gmail.com}
\homepage{https://www.lucamoretti.ai}

\address{Via Alessandro Volta 42, Milano, Lombardia, Italia, 20154}{}{}

\extrainfo{{\small \faLinkedin}\ \href{https://www.linkedin.com/in/lucamoretti-ai}{@lucamoretti-ai} {} {} {} • {} {} {} 
{\small \faGithub}\ \href{https://github.com/lucamoretti}{@lucamoretti}}

\begin{document}

\maketitle

\section{Informazioni personali}

\cvline{}{Ingegnere AI con oltre 7 anni di esperienza nella progettazione e nel deployment di sistemi di apprendimento automatico scalabili. Specializzato in modelli di linguaggio, visione artificiale e pipeline MLOps in ambienti cloud.
%
\begin{itemize}
\item Guida tecnica di team distribuiti e mentorship di giovani ingegneri.
\item Forte orientamento ai risultati, alla qualità del codice e all'impatto di business.
\item Esperienza end-to-end: dalla raccolta dati alla messa in produzione e monitoraggio.
\end{itemize}}
\section{Formazione}

\cventry{set 2017–lug 2019}
        {Laurea magistrale, Ingegneria Informatica, Punteggio: 110/110 con lode}
        {Politecnico di Milano}
        {\url{https://www.polimi.it}}
        {}
        {\begin{itemize}
\item Tesi magistrale su modelli transformer per il rilevamento di anomalie in serie temporali industriali.
\item Progetto di gruppo con azienda manifatturiera per la manutenzione predittiva.
\item Media ponderata elevata e attività di tutoraggio per corsi di programmazione.
\end{itemize}
\textbf{Corsi}: Apprendimento automatico e riconoscimento di pattern, Elaborazione del linguaggio naturale, Visione artificiale, Sistemi distribuiti e cloud computing, Ottimizzazione e ricerca operativa, Architetture dei calcolatori}

\cventry{set 2014–lug 2017}
        {Laurea triennale, Ingegneria Informatica, Punteggio: 108/110}
        {Università di Bologna}
        {\url{https://www.unibo.it}}
        {}
        {\begin{itemize}
\item Solida base in matematica applicata, algoritmi e sviluppo software.
\item Progetto finale su sistemi di raccomandazione basati su filtraggio collaborativo.
\end{itemize}
\textbf{Corsi}: Algoritmi e strutture dati, Basi di dati, Reti di calcolatori, Probabilità e statistica, Programmazione orientata agli oggetti}
\section{Esperienza professionale}

\cventry{gen 2022–Presente}
        {Ingegnere AI Senior}
        {Neuralis Labs}
        {\url{https://www.neuralislabs.it}}
        {}
        {\begin{itemize}
\item Progettazione e implementazione di una piattaforma di modelli di linguaggio per l'assistenza clienti, con riduzione del 35\% dei tempi di risposta.
\item Responsabile tecnico di un team di 6 ingegneri AI e data scientist, con mentorship e revisione del codice.
\item Definizione di architetture MLOps su Kubernetes e AWS per training, deployment e monitoraggio continuo.
\item Collaborazione con product manager e stakeholder per tradurre requisiti di business in soluzioni AI misurabili.
\end{itemize}
\textbf{Parole chiave}: LLM, MLOps, Kubernetes, AWS, NLP}

\cventry{mar 2020–dic 2021}
        {Ingegnere Machine Learning}
        {DataMind Italia}
        {\url{https://www.datamind.it}}
        {}
        {\begin{itemize}
\item Sviluppo di modelli di visione artificiale per il controllo qualità in linea di produzione, con accuratezza superiore al 96\%.
\item Creazione di pipeline di dati automatiche con Python, Airflow e Spark su Azure.
\item Implementazione di test A/B e monitoraggio delle performance dei modelli in produzione.
\item Documentazione tecnica e trasferimento di conoscenza ai team operativi.
\end{itemize}
\textbf{Parole chiave}: Computer Vision, Python, Apache Airflow, Azure, Test A/B}

\cventry{set 2019–feb 2020}
        {Scienziato dei Dati}
        {Aurora Robotics}
        {\url{https://www.aurorarobotics.it}}
        {}
        {\begin{itemize}
\item Analisi di dati sensoriali per robot collaborativi e sviluppo di modelli predittivi per la manutenzione.
\item Prototipazione di algoritmi di rilevamento anomalie con scikit-learn e PyTorch.
\item Presentazione dei risultati a team tecnici e non tecnici.
\end{itemize}
\textbf{Parole chiave}: Data Science, PyTorch, Scikit-learn, Anomaly Detection}
\section{Lingue}

\cvline{Italiano}{Madrelingua o bilingue \hfill \textbf{Parole chiave}: Madrelingua}
\cvline{Inglese}{Competenza professionale completa \hfill \textbf{Parole chiave}: IELTS 8.0, C1}
\cvline{Francese}{Competenza professionale limitata \hfill \textbf{Parole chiave}: DELF B1}
\section{Competenze}

\cvline{Apprendimento automatico}{Esperto \hfill \textbf{Parole chiave}: PyTorch, TensorFlow, Scikit-learn, XGBoost, Reti neurali}
\cvline{MLOps}{Avanzato \hfill \textbf{Parole chiave}: Docker, Kubernetes, MLflow, Kubeflow, CI/CD}
\cvline{Elaborazione del linguaggio naturale}{Avanzato \hfill \textbf{Parole chiave}: Transformer, LLM, Hugging Face, RAG, Fine-tuning}
\cvline{Cloud e infrastrutture}{Avanzato \hfill \textbf{Parole chiave}: AWS, Azure, GCP, Terraform, Serverless}
\cvline{Data Engineering}{Intermedio \hfill \textbf{Parole chiave}: SQL, Spark, Airflow, Kafka, Snowflake}
\section{Premi}

\cventry{nov 2023}
        {Associazione Italiana per l'Intelligenza Artificiale}
        {Premio Innovazione AI}
        {}
        {}
        {Riconoscimento per il progetto di assistenza clienti basato su modelli linguistici, selezionato tra oltre 120 candidature nazionali.}
\section{Certificazioni}

\cventry{feb 2023}
        {Amazon Web Services}
        {AWS Certified Machine Learning - Specialty}
        {\url{https://aws.amazon.com/certification/}}
        {}
        {}

\cventry{giu 2021}
        {Google}
        {TensorFlow Developer Certificate}
        {\url{https://www.tensorflow.org/certificate}}
        {}
        {}
\section{Pubblicazioni}

\cventry{set 2022}
        {Modelli transformer per il rilevamento di anomalie industriali}
        {Atti del Congresso Italiano di Intelligenza Artificiale}
        {\url{https://doi.org/10.1000/aiia.2022.042}}
        {}
        {\begin{itemize}
\item Studio su architetture transformer applicate a serie temporali multivariate per la manutenzione predittiva.
\item Confronto con modelli tradizionali e analisi dell'impatto della qualità dei dati.
\item Dataset e codice resi disponibili alla comunità scientifica.
\end{itemize}}
\section{Progetti}

\cventry{apr 2021–nov 2021}
        {Sistema di controllo qualità visivo basato su deep learning per linee di produzione industriali}
        {VisionQC}
        {\url{https://github.com/lucamoretti/visionqc}}
        {}
        {\begin{itemize}
\item Rilevamento automatico di difetti superficiali con reti convoluzionali e segmentazione semantica.
\item Integrazione con telecamere industriali e dashboard di monitoraggio in tempo reale.
\item Riduzione dei falsi negativi del 28\% rispetto al sistema precedente.
\end{itemize}
\textbf{Parole chiave}: Computer Vision, Deep Learning, PyTorch, FastAPI}

\cventry{feb 2023–ago 2023}
        {Assistente conversazionale con RAG per la ricerca semantica su documenti aziendali}
        {Assistente Documentale}
        {\url{https://github.com/lucamoretti/doc-assistant}}
        {}
        {\begin{itemize}
\item Pipeline di ingestion, chunking e indicizzazione vettoriale di documenti PDF e testuali.
\item Integrazione di modelli LLM con controllo delle citazioni e mitigazione delle allucinazioni.
\item Interfaccia web e API per l'uso da parte dei team interni.
\end{itemize}
\textbf{Parole chiave}: RAG, LLM, Vector Database, LangChain, Python}
\section{Interessi}

\cvline{Tecnologia}{Open source, Robotica, Intelligenza artificiale}
\cvline{Sport}{Corsa, Ciclismo, Arrampicata}
\cvline{Musica}{Chitarra, Jazz}
\section{Volontariato}

\cventry{ott 2021–Presente}
        {Mentore AI}
        {PoliAi}
        {\url{https://www.poliai.it}}
        {}
        {\begin{itemize}
\item Mentorship per studenti universitari su progetti di machine learning e carriera nel settore AI.
\item Organizzazione di workshop su Python, PyTorch e buone pratiche MLOps.
\item Supporto alla community con revisione di progetti open source.
\end{itemize}}
    
\end{document}